% !TeX root = ../main.tex
\chapter{Introduction}
\label{chap:introduction}

La mesure de température est courante et critique dans les procédés industriels, les bancs d'essais et les systèmes électroniques. Ce laboratoire est consacré à l'étude pratique et à la comparaison de deux technologies de capteurs thermorésistifs usuelles :
\begin{itemize}[leftmargin=1.2cm]
    \item Une \textbf{thermistance CTN} (perle miniature SEMI833-ET) : capteur semi-conducteur à coefficient de température négatif, présentant une forte sensibilité mais une réponse très non-linéaire.
    \item Une \textbf{sonde RTD Pt100} (modèle industriel KP chemisé en acier inoxydable) : résistance de platine à coefficient positif, réputée pour sa linéarité, sa reproductibilité et sa stabilité dans le temps (norme IEC 60751).
\end{itemize}

Les manipulations sont réalisées avec le matériel du laboratoire comprenant un multimètre de table Keysight 34460A (résolution 6½ digits), un four d'étalonnage à sec Fluke 9103 et une carte d'acquisition National Instruments PCIe-6321 interfacée sous LabVIEW.

\section{Objectifs du travail}
\label{sec:intro_objectifs}

Le travail pratique poursuit six objectifs principaux :
\begin{enumerate}[leftmargin=1.2cm]
    \item \textbf{Mesure ambiante par la CTN :} Mesurer manuellement la résistance de la thermistance à température ambiante et déduire la température par la pente locale issue de la datasheet constructeur.
    \item \textbf{Conditionnement et auto-échauffement :} Dimensionner un pont diviseur sous \qty{14}{\volt} pour adapter le signal à la carte d'acquisition et vérifier par le calcul que l'auto-échauffement interne par effet Joule reste négligeable devant la tolérance du capteur.
    \item \textbf{Acquisition LabVIEW et suréchantillonnage :} Mettre en œuvre la chaîne d'acquisition avec conversion par table de correspondance, puis comparer une mesure normale et une mesure avec suréchantillonnage (moyenne de 100 points) lors du chauffage au sèche-cheveux.
    \item \textbf{Étalonnage au point fixe à \qty{0}{\degreeCelsius} de la Pt100 :} Calibrer la sonde dans le puits sec pour déterminer sa résistance réelle $R_0$.
    \item \textbf{Comparaison métrologique 2 fils vs 4 fils :} Mesurer l'erreur systématique introduite par les conducteurs de liaison en montage deux fils (2W) et justifier l'intérêt du montage quatre fils (4W) Kelvin.
    \item \textbf{Caractérisation dynamique de la Pt100 :} Mesurer la réponse indicielle lors d'un échelon thermique en eau chaude et extraire la constante de temps $\tau$ ainsi que le temps de réponse à \qty{95}{\percent} ($t_{95\%}$) par différentes méthodes (tangente, seuil à \qty{63.2}{\percent}, ajustement par moindres carrés).
\end{enumerate}
